\chapter{Experimental discovery and the route to proof}\label{ch:10-discovery}
The mathematical progression of the manuscript suggests a practical method.
One begins with a value that is easy to compute but hard to recognize, finds
an exact transformation that organizes its arguments, and then asks which
coordinates remain after known laws are imposed. Numerical experiments are
especially effective at proposing those coordinates. Their usefulness
increases when the proposed identity is subsequently derived, and when a
negative result is recorded with enough detail to explain its limits.

\section{Constructing a meaningful integer-relation experiment}
An experiment consists of a target, a finite ordered basket of constants,
their coefficient field, a working precision, a coefficient bound and an
acceptance rule. Weight and character parity give strong guidance for the
basket. For instance a rational-angle Fourier transform can require
$\sqrt q$-weighted atoms; searching only with rational coefficients in the
unweighted values misses that structure. The
$\pi^3\psi^{(-3)}(1/4)$ coordinate in Chapter~\ref{ch:08-differentiation}
is another example: omitting the $\pi^3$ factor changes the question.

Before searching, remove known basket dependencies. A search containing
both $\log\Gamma(1/2)$ and $\frac12\log\pi$ may return their tautology
instead of a relation involving the target. Reflecting arguments,
distributing roots of unity and replacing elementary odd beta values
reduce this problem before numerical linear algebra starts.

Working precision must exceed the cancellation loss as well as the digits
needed to discriminate the proposed coefficient height. The heuristic
$10^{P/d}$, for $P$ available digits and $d$ basket coordinates, estimates
a scale at which accidental short numerical relations can appear. It is
neither a theorem about a specific basket nor a rigorous exclusion bound.
Coefficient sizes, conditioning and the number of searches all matter.
The earlier false negative for $S_5$ illustrates this: a $51$-coordinate
basket at $150$ digits obscured the height-$3840$ relation that a smaller,
properly structured basket exposed.

\section{Controls and independent recomputation}
A planted positive control checks that the search can recover a known
relation of comparable scale. A perturbation of one coefficient tests the
resolution of the residual check. A negative control and an auxiliary
``canary'' constant can expose numerical coincidences. These controls are
diagnostics. In particular no independence theorem is known between the
Champernowne constant and the periods in these baskets, and a zero canary
coefficient cannot certify an identity.

A successful candidate should be frozen as exact rational or algebraic
coefficients, then reevaluated at higher precision using an independent
representation. For a nested sum that may be an iterated integral; for a
CM row sum it may be a modular Fourier series; for Herglotz values it may
be a defining integral compared with the finite dilogarithm sum.
Agreement between two engines using the same incorrect formula is less
informative than agreement between different mathematical representations.
Error-controlled interval evaluation can establish a finite residual
bound, but even an arbitrarily small nonzero numerical bound alone does
not establish equality of general transcendental constants.

\section{What exact linear algebra proves}
An exact relation matrix proves an upper bound on the number of necessary
coordinates in its presented quotient. Its row-combination matrix is a
replayable finite certificate relative to the input identities. A free
column is independent in that formal quotient, not automatically among
the evaluated numerical periods. Additional functional equations or
regularization can reduce the quotient further.

Likewise a motivic dimension theorem concerns motivic objects. To infer
numerical independence one needs a period-injectivity statement, generally
conjectural. A dimension for a chosen word-length grading need not be the
dimension of the classical depth filtration. The Gaussian triples are a
natural test case for this distinction.

Certificates require both exact atom cancellation and sound law
instantiation. Substitution strings must be parsed within a restricted
grammar; branches and domains must be established mathematically. A
numerical branch guard is useful for rejecting errors but leaves a
separate domain obligation. These issues motivate a future formal
development rather than a current claim of proof-assistant verification.

\section{Computational conventions}
The manuscript uses decreasing nested indices. In the tested Wolfram
Language~15 interface,
\[
\Li_{a,b}(z_1,z_2)=
\texttt{MultiplePolyLog[\{a,b\},\{z1,z2\}]},\qquad
\zeta(a,b)=\texttt{MultipleZeta[\{a,b\}]}.
\]
Finite checks such as $\zeta(2,1,1)=\zeta(4)$ and
$H(0,1;z)=\Li_2(z)$ calibrate the interface. A divergent call does not
resolve ordering ambiguities on its own. The same calibration is needed
for signed harmonic indices and generalized-polylogarithm letter order.

Negative-order polygamma means the integral based at zero:
$\psi^{(-1)}=\log\Gamma$ and
$\psi^{(-n)}(z)=\int_0^z\psi^{(-(n-1))}(t)\,dt$ for $n\ge2$.
Balanced conventions can differ by explicit polynomials. In Hurwitz jets,
$\zeta^{(j)}$ always differentiates the spectral variable; the superscript
$(k)$ on $\gamma_n$ differentiates its parameter.
The CM hexagonal point $e^{i\pi/3}$ and the Eisenstein cube root
$e^{2\pi i/3}$ are different arguments even when a local section denotes
each by $\rho$.

\section{The surviving research programme}
Several concrete proof problems now fit into one framework. The odd-weight
$S_p$ directions need a precise choice of comparison algebra and a
rigorous reduction or obstruction; the even-weight direction has the
generating-integral proof of Chapter~\ref{ch:04-depth}. The Gaussian and
Eisenstein mixed points need a full regularized relation computation with
an explicit matrix and certified reductions. The Stieltjes distribution
ranks need a uniform proof of the divisor-row structure. Numerical
independence of character derivative coordinates remains a separate
arithmetic problem.

The gamma programme can transport its exact row certificates across the
DFT and prove branch validity for each law instance. Its finite-matrix
termination should remain separate from Rohrlich's conjecture. The CM
programme can replace fitted coefficients by modular equations and
class-field computations, while taking account of periods under Galois
conjugation. The Herglotz programme needs complete explicit class-character
and unit definitions for the regulator candidates before they become
locally replayable, followed by the relevant Stark argument.

The cubic log-gamma moment and the opaque individual CM components are
good discovery targets because their defining integrals are accessible,
but the naive single-value baskets miss their likely depth or modular
structure. A richer representation is often more productive than a
larger unstructured search. The scientific value of the experiments is
that they tell us which representation to seek next.
